\section{Imbalance two: returns of span at most two}\label{sec:spantwo}
In this section, we prove that every orientation $++W$ in the family~\ref{ov:family1}, where $W$ has span at most two, has a perfect fractional homomorphic tiling in every digraph $F$ with $\delta^0(F)>|F|/2$ (Theorem~\ref{height390:two-width-factor}).
For the four short orientations $Q=+++-$, $++-++-$, $D_{4,2}$ and $D_{5,3}$ we argue as for directed cycles, and the longer orientations then follow either by Lemma~\ref{transfer:fibres} or from the bounds of Theorem~\ref{norm391:normalized-cover}.

\subsection{The four-cycle \texorpdfstring{$Q$}{Q}}\label{q:section}
Recall that $Q$ has edges $ab,bc,cd,ad$. Deleting a position from $Q$ leaves a path with two edges, whose middle vertex is a source, a sink, or internal. So a partial homomorphism of $Q$ is a triple $xyz$ of vertices in which $x,y$ and $y,z$ are adjacent in a suitable direction, where $x=z$ is permitted; we call such a triple a \emph{three-position walk}.
If $x\to y\to z$, there are two partial homomorphisms with this image, with deleted position $a$ or $d$, and their completion pairs are $(N^-(x),N^-(z))$ and $(N^+(x),N^+(z))$.
If $x\to y\leftarrow z$ or $x\leftarrow y\to z$, the completion pairs are $(N^+(x),N^-(z))$ and $(N^-(x),N^+(z))$.

\begin{lemma}\label{q:endpoint}
Let $D$ be a digraph with nonnegative costs $c$. Suppose that every three-position walk
costs at least $s$, every $2$-cycle costs more than $s$, and every
homomorphism from $Q$ costs more than $2s$. Let $\rho$ be a nonzero
measure of mass $h$ on three-position walks of cost $s$, with multiplicity function $p$. Then some three-position walk of cost $s$ has a completion pair $(X,Y)$ with
$p(X)+p(Y)\le2h$.
\end{lemma}

\begin{proof}
Restrict $D$ to the support of $p$, so that $0\le c(v)\le s$ for every $v$ and $p(V(D))=3h$.
Every completion pair of a three-position walk of cost $s$ is disjoint, as a common
vertex would give a homomorphism from $Q$ of cost at most $2s$.
Suppose that every such pair has total multiplicity greater than $2h$.

Apply Lemma~\ref{norm391:cutoff} in case~\ref{norm391:cutoff:1}, where each three-position walk is regarded as a partial homomorphism with any of its deleted positions. Its proof subtracts some $\alpha\ge0$ from all costs, and so $3\alpha$ from $s$, and the strict lower bound on the costs of $2$-cycles is preserved. Let $g>0$ be the mass of the new measure and $s'$ the new minimum cost.
Every three-position walk in the support of the new measure meets the set $Z$ of vertices of cost zero.

Suppose first that $s'>0$. If $z\in Z$ and $u$ is a neighbour of $z$, then the walk $zuz$ shows that $c(u)\ge s'$. Thus every walk in the new support has costs $(0,s',0)$ in order: a walk of cost $s'$ with its middle vertex in $Z$ would cost at least $2s'$.
So the new support consists of a set $Z$ of multiplicity $2g$ and a set of vertices of cost $s'$ of multiplicity $g$. There are no edges inside $Z$, as a walk $zz'z$ would cost zero, and no $2$-cycles between the two sets, as $2$-cycles cost more than $s'$.
For adjacent $z\in Z$ and $u$, the walk $zuz$ has completion pair $(N^+(z),N^-(z))$, which is disjoint from $Z$, and so has total multiplicity at most $g$, a contradiction.

Suppose now that $s'=0$, so that all costs are zero. Then the new support is an oriented graph and receives no homomorphism from $Q$.
By Lemma~\ref{general:tournament-paths} with $q=1$, its underlying graph is $K_4$-free.
Lemma~\ref{general:weighted-turan}, with total multiplicity $3g$, gives
a vertex $z$ of positive multiplicity with
$p(N^+(z))+p(N^-(z))\le2g$.
As $z$ lies on a walk in the support, it has a neighbour $u$, and the walk $zuz$ contradicts our assumption.
\end{proof}

\begin{theorem}\label{q:fractional}
Every digraph $F$ of order $k$ with $\delta^0(F)>k/2$ has a perfect
fractional homomorphic $Q$-tiling.
\end{theorem}

\begin{proof}
Suppose not. By Lemma~\ref{general:LP}, there is a deficient optimal cover $w$, with $0\le w\le1$, total weight
$W<k/4$ and $\min w=0$.
Let $M=\max w$ and $s=1-M$. Every three-position walk can be completed
to a homomorphism from $Q$, because each of its completion sets has more than $k/2$ vertices, so its cost is at least $s$.
Moreover, $M>1/2$. Indeed, for a vertex $o$ of weight zero, every
$u\in N^+(o)$ has $w(u)\ge s$, by the walk $ouo$, so, if $M\le1/2$, then
$W\ge |N^+(o)|/2>k/4$.

Let $C$ be the set of vertices of weight $M$ and $h=|C|$. As $2M>1$, every tight homomorphism from $Q$
has at most one position in $C$. Let $\rho$ be the measure of
Lemma~\ref{general:max-deletion}, which has mass $h$ and is supported on three-position walks of cost $s$ in $F-C$.
We check the hypotheses of Lemma~\ref{q:endpoint} in $F-C$. A $2$-cycle
$uv$ receives the homomorphism $(a,b,c,d)\mapsto(u,v,u,v)$ from $Q$, so
$w(u)+w(v)\ge1/2>s$, and every homomorphism from $Q$ costs at least
$1>2s$. Thus Lemma~\ref{q:endpoint} gives a three-position walk of cost $s$ with a completion pair of total multiplicity at most $2h$, contradicting
Lemma~\ref{general:max-deletion}.
\end{proof}


\subsection{Three short orientations}\label{sec:threemixed}
Let
\[
 B=++-++-,\qquad A=D_{4,2}=++++--,\qquad D=D_{5,3}=+++++---.
\]
We use the partial homomorphisms and completion sets of
Section~\ref{norm391:section}, and treat these three orientations together.
We use the following partial homomorphisms. Delete the first vertex of $B$. Reading the remaining path from the next vertex, and translating heights by two after we wrap around, its heights are $1,0,1,2,1$. For $A$, delete the $j$th vertex, where $j\in\{1,2,3\}$; the remaining path climbs from height $j-1$ to $2$, comes down to $0$ and climbs back to $j-1$ (for example, its heights are $1,2,1,0,1$ when $j=2$). For $D$, we do the same with $j\in\{1,2,3,4\}$, and the heights now range over $0,1,2,3$.
In each case both ends of the path are at the same height, and every edge of $H$ goes from a vertex at some height to a vertex at the next height. So, sending each height to the corresponding vertex of a directed walk with two edges (for $B$ and $A$) or three edges (for $D$), we obtain a partial homomorphism into this walk whose two endpoints have the same image, $z$ say, and whose completion sets are $N^-(z)$ and $N^+(z)$. Every vertex of the walk other than its two ends can play the role of $z$.
In particular, for each of these orientations, the span of $H-j$ is at most three for some $j$.

\begin{lemma}\label{base392:zero}
Let $H\in\{B,A,D\}$ have order $r$, and let $D_0$ be a digraph which receives no homomorphism from
$H$. Then, for every nonzero measure of mass $m$ on partial homomorphisms into $D_0$, some partial homomorphism $Q$ has $\Phi_p(Q)\le(r-2)m$.
\end{lemma}
\begin{proof}
Restrict $D_0$ to the support of $p$, and suppose that every partial homomorphism has $\Phi_p$-value greater than $(r-2)m$. The digraph $D_0$ is oriented, since every even cycle maps to a $2$-cycle. For $z\in V(D_0)$, write $d_p(z)=p(N^-(z))+p(N^+(z))$.
The underlying graph of $D_0$ is $K_4$-free if $H=B$, $K_5$-free if $H=A$, and $K_6$-free if $H=D$. Indeed, $B$ maps to a transitive triangle by $(a,b,c,a,b,c)$, and every tournament on four vertices contains a transitive triangle, as some vertex has two inneighbours or two outneighbours. Moreover, Lemma~\ref{general:tournament-paths}
supplies homomorphisms from $D_{4,2}$ and $D_{5,3}$ to tournaments on five and six vertices.

\emph{The case $H=B$.} Here $P=5m$. Let $E$ be the set of vertices
with an inneighbour and an outneighbour. For $z\in E$, the partial homomorphism of the table with both endpoints at $z$ gives $d_p(z)>4m$.
By Lemma~\ref{base392:extremes}, the vertices outside $E$ have multiplicity less than $2m$, so $p(E)>3m$. For $z\in E$, the multiplicity of the neighbourhood of $z$ inside $E$ is greater than $4m-(P-p(E))=p(E)-m$. However, Lemma~\ref{general:weighted-turan} gives a vertex $z\in E$ whose neighbourhood inside $E$ has multiplicity at most $2p(E)/3$. Thus $p(E)<3m$, a contradiction.

\emph{The case $H=A$.} Suppose first that every vertex lies on a directed walk with two edges. Then, by the table, every vertex $z$ has $d_p(z)>4m$. However, Lemma~\ref{general:weighted-turan} gives a vertex with $d_p(z)\le 3P/4=15m/4$, a contradiction.
So some vertex $f$ lies on no directed walk with two edges, and $f$ is the image of a position of some partial homomorphism in the support. The sign words of the partial homomorphisms of $A$ are $++++$, $+++-$, $++--$, $+--+$, $--++$ and $-+++$. In the words $++++$, $++--$ and $--++$, every position lies on a directed path with two edges. So the position mapped to $f$ is an end position of a partial homomorphism with word $+++-$, $+--+$ or $-+++$, and it is an endpoint adjacent to the deleted position. The completion set at this endpoint is the inneighbourhood or the outneighbourhood of $f$ which is empty, as $f$ lies on no directed walk with two edges.
On the other hand, for any edge $v\to u$, no partial homomorphism in the support lies entirely in $N^-(v)$. Indeed, every partial homomorphism contains a directed path $x\to y\to z$, and then the two paths $x\to y\to z\to v\to u$ and $x\to v\to u$ would give a homomorphism from $A=D_{4,2}$. Thus $p(N^-(v))\le4m$, and, similarly, $p(N^+(u))\le 4m$. In the partial homomorphism above, the other completion set is of this form, so its $\Phi_p$-value is at most $4m$, a contradiction.

\emph{The case $H=D$.} Here $P=7m$. Let $E$ be the set of vertices
lying on a directed walk with three edges, and $F_0=V(D_0)\setminus E$.
Every partial homomorphism of $D$ contains a directed path with three edges, so $E\ne\varnothing$.
By the table, every $z\in E$ has $d_p(z)>6m$, that is, $P-d_p(z)<m$. Fix a walk $z_0\to z_1\to z_2\to z_3$.
For $f\in F_0$, let $a(f)$ and $c(f)$ be the maximum lengths of directed walks
ending and starting at $f$, respectively. These are finite and satisfy
$a(f)+c(f)\le2$, as otherwise concatenating two such walks would put $f$ in $E$.
An edge $z_i\to f$ requires $i+1\le a(f)$, and an edge
$f\to z_i$ requires $4-i\le c(f)$. Hence $f$ is not adjacent
to at least two of $z_0,z_1,z_2,z_3$, even if some of them coincide. As $D_0$ is oriented, double counting gives
\[
 2p(F_0)\le\sum_{i=0}^3(P-d_p(z_i))<4m,
 \qquad\text{so}\qquad p(E)>5m.
\]
Every $z\in E$ thus has more than $p(E)-m$ multiplicity in its neighbourhood inside $E$, whereas
Lemma~\ref{general:weighted-turan} gives a vertex $z\in E$ with at most $4p(E)/5$. Hence
$p(E)<5m$, a contradiction.
\end{proof}

\begin{lemma}
\label{base392:endpoint}
\label{b:endpoint}\label{a:endpoint}\label{repair:d53-endpoint}
Let $H\in\{B,A,D\}$ have order $r$, and let $D$ be a digraph with nonnegative costs such that every partial homomorphism costs at least $L\ge0$ and every
homomorphism from $H$ costs more than $2L$.
Then, for every nonzero measure of mass $m$ on partial homomorphisms of cost $L$, some
minimum partial homomorphism $Q$ has $\Phi_p(Q)\le(r-2)m$.
\end{lemma}
\begin{proof}
Suppose not, and apply Lemma~\ref{norm391:cutoff} in case~\ref{norm391:cutoff:1}. We obtain a measure of positive mass, supported on minimum partial homomorphisms which meet the set $Z$ of vertices of cost zero, with new minimum cost $L'\ge0$, such that every homomorphism from $H$ costs more than $2L'$ and~\eqref{norm391:counter} holds.
If $L'>0$, then, as the span of $H-j$ is at most three for some $j$, Lemma~\ref{norm391:positive-endpoint} in case~\ref{norm391:positive-endpoint:1} gives a contradiction.
If $L'=0$, then all costs in the support are zero, the support receives no homomorphism from $H$, and Lemma~\ref{base392:zero} gives a contradiction.
\end{proof}

\begin{lemma}
\label{base392:maximum}\label{aux:d53maximum}
Let $H\in\{B,A,D\}$, and let $F$ be a digraph with $\delta^0(F)>|F|/2$. Then every deficient optimal fractional $H$-cover of $F$
has $M=\max w>1/2$. If $H\in\{B,A\}$, then $M>2/3$.
\end{lemma}
\begin{proof}
By Lemma~\ref{general:LP}, we have $0\le w\le1$ and $\min w=0$. Every partial homomorphism costs at least $L=1-M$, since
its two completion sets intersect. Let $d=\delta^0(F)$ and
$\tau=(\sum w)/d<2/r$, and let $o$ be a vertex of weight zero, with minimum weights $\alpha$ and $\beta$ in $N^-(o)$ and $N^+(o)$. We use Lemma~\ref{norm391:savings}.

For $B$ and $A$, the partial homomorphism of the table which maps the endpoints to $o$, and the other two heights to minimum-weight neighbours of $o$, costs
$\alpha+\beta$. Hence $L\le\alpha+\beta<\tau<1/3$, so $M>2/3$.

For $D$, reverse the edges of $F$ if necessary so that $\alpha\le\beta$;
this is possible as $D_{5,3}$ is isomorphic to its converse. Let $z$ be
a vertex of minimum weight $a$ in $N^-(o)\cap N^+(o)$.
Choose $u\to o\to v$ with $w(u)=\alpha$ and $w(v)=\beta$. For each $x\in N^-(u)$, the two paths
\[
 x\to u\to o\to z\to o\to v\qquad\text{and}\qquad x\to u\to o\to v
\]
give a homomorphism from $D_{5,3}$, and so $w(x)\ge1-a-2\alpha-\beta$. Thus the minimum weight in $N^-(u)$, which is at most $\chi$, is at least $1-a-2\alpha-\beta$.
By Lemma~\ref{norm391:savings}, $\chi+\alpha<\tau$ and $\alpha+\beta<\tau$. Therefore $a>1-2\tau>1/2$, and so $M>1/2$.
\end{proof}

\begin{theorem}
\label{base392:fractional}
\label{b:fractional}\label{a:fractional}\label{repair:d53-factor}
Every digraph $F$ with $\delta^0(F)>|F|/2$ has perfect fractional homomorphic
tilings by $B=++-++-$, by $D_{4,2}$ and by $D_{5,3}$.
\end{theorem}
\begin{proof}
Let $H\in\{B,A,D\}$ have order $r$, and suppose that $F$ has no perfect fractional homomorphic $H$-tiling. By Lemma~\ref{general:LP}, there is a deficient optimal cover $w$, and Lemma~\ref{base392:maximum} gives $1/2<M\le1$.
Let $C$ be the set of vertices of weight $M$ and $h=|C|$. Since $2M>1$, every tight homomorphism has at most one position in $C$, and so the measure of
Lemma~\ref{general:max-deletion} has mass $h$ and is supported on
partial homomorphisms of cost $L=1-M$ in $F-C$, which is the minimum cost there.
Every homomorphism from $H$ costs at least $1>2L$, so
Lemma~\ref{base392:endpoint} gives a minimum partial homomorphism with $\Phi_p$-value at most $(r-2)h$, which contradicts
Lemma~\ref{general:max-deletion}.
\end{proof}


\subsection{Isolated backward edges}
\label{transfers}

\begin{theorem}\label{imbalance402:isolated}
\label{transfer:Qfamily}\label{transfer:Bfamily}
Let $H$ be an orientation of an even cycle with $k_H=2$, and suppose that, in a reading with sign sum two, no two backward edges are consecutive. Then every digraph $F$ with $\delta^0(F)>|F|/2$ has a
perfect fractional homomorphic $H$-tiling.
\end{theorem}
\begin{proof}
If $H$ has $t$ backward edges, then it has $t+2$ forward edges, which form $t$ directed runs of length at least one, separated by the isolated backward edges. Thus either one run has length three, or two runs have length two, and all
other runs have length one. Up to rotation, $H$ is therefore one of
\[
 H_j=+++(-+)^j-\qquad\text{and}\qquad
 T_{a,b}=++(-+)^a-++(-+)^b-,\qquad\text{where }j,a,b\ge0.
\]
We show that perfect fractional tilings by $Q$ and $B$ give perfect fractional tilings by $H_j$ and $T_{a,b}$, respectively, using Lemma~\ref{transfer:fibres}. As Theorems~\ref{q:fractional} and~\ref{b:fractional} provide these tilings, this proves the theorem.

A long alternating path can be mapped onto a single edge, as follows. Suppose that $u\to v$ is an edge of a digraph $K$, and that we replace it by a path of length $2m+1$ whose edges alternate in direction, starting and ending with forward edges. If $u\to z$ is an edge of $K$, we can map this path to $u,(z,u)^m,v$, where $(z,u)^m$ denotes $m$ repetitions of $z,u$; and if $z\to v$ is an edge of $K$, we can map it to $u,v,(z,v)^m$. In both cases, the new path places $m$ further positions at each end of the edge between $u$ or $v$ and $z$.

For $H_j$, replace the edge $a\to d$ of
$Q=(a\to b\to c\to d,\ a\to d)$ by such a path of length $2j+1$.
Map it using the edge $ab$ or the edge $cd$, each with probability one half.
Each vertex of $Q$ then receives $j/2$ further positions in expectation, so its expected fibre size is $1+j/2=|H_j|/4$.

For $T_{a,b}$, write
$B=0\to1\to2\leftarrow3\to4\to5\leftarrow0$, and let $\sigma(i)=i+3\pmod6$, which is an automorphism of $B$.
Replace the edges $3\to2$ and $0\to5$ by alternating paths of lengths $2a+1$ and $2b+1$, respectively.
Choose $\gamma$ uniformly from $\{\mathrm{id},\sigma\}$ and an edge $f$
uniformly from $\{12,32,34\}$, each of which meets $3\to2$.
Map the old positions by $\gamma$, and map the two new paths using the edges
$\gamma(f)$ and $\gamma\sigma(f)$, respectively.
Each of these two edges is uniformly distributed over the six edges of $B$, so each vertex of $B$ is an end of each of them with probability
$1/3$. Thus each vertex has expected fibre size
$1+(a+b)/3=|T_{a,b}|/6$.

Finally, if the reading of $H$ with sign sum two is its converse reading, we reverse the edges of both $H$ and $F$. Both $Q$ and $B$ are isomorphic to their converses, via $(a\ d)(b\ c)$ and
$i\mapsto5-i\pmod6$, respectively, so the same argument applies.
\end{proof}

\subsection{Returns of span two}

The remaining orientations in~\ref{ov:family1} with span at most two are handled using Theorem~\ref{norm391:normalized-cover}.
Let $H$ be a cycle orientation with sign sum two, and let $j$ be a position of $H$. For each height $i$ of the height walk of $H-j$, the \emph{height class} of $i$ is the set of positions at height $i$, and its multiplicity is its size. If the height walk of $H-j$ crosses between heights $i$ and $i+1$ exactly $e_i$ times in each direction, then the class of $i$ has multiplicity $e_{i-1}+e_i$, plus one if $i$ is the height of the last position.

\begin{lemma}
\label{height390:heavy-pair}
Let $H$ have order $2S$ and sign sum two, and satisfy
\eqref{norm391:range}. Suppose that, for some position $j$, two adjacent height classes of $H-j$ have multiplicities $p$ and $q$ with
\[
 p,q\ge S/2\qquad\text{and}\qquad p+q\ge3S/2-1.
\]
Then every digraph $F$ with $\delta^0(F)>|F|/2$ has a perfect fractional homomorphic
$H$-tiling. The same holds if, instead, a convex combination of the vectors of class multiplicities of several partial homomorphisms satisfies these inequalities.
\end{lemma}
\begin{proof}
Suppose not. By Lemma~\ref{general:LP}, there is a deficient optimal cover $w$. Let
$\tau=(\sum w)/\delta^0(F)<1/S$, $M=\max w$ and
$\kappa=\min_{xy\in E(F)}(w(x)+w(y))$.
Theorem~\ref{norm391:normalized-cover} and
Lemma~\ref{norm391:savings} give
\[
 2M+(S-2)\kappa\le1,\qquad
 \chi+\kappa<\tau,\qquad 2\kappa<\tau.
\]
Map the two adjacent height classes to the two ends of an edge of weight $\kappa$. Then map the other heights in turn to neighbours of weight at most $\chi$, so that each edge of $H-j$ goes from one height to the next, and complete the deleted position by a vertex of weight
at most $M$. With $A=\max(p,q)$ and $C=2S-1-p-q$, the cost of this homomorphism shows that $1\le M+A\kappa+C\chi$. This also holds for a convex combination of class multiplicity vectors, by averaging.
The hypotheses give $0\le C\le S/2$ and
$A+C=2S-1-\min(p,q)\le3S/2-1$. Therefore,
\[
\begin{split}
 1&\le(2A-S+2)\kappa+2C\chi\\
  &\le2C(\chi+\kappa)+2(S-2C)\kappa
   <S\tau<1,
\end{split}
\]
where the strict inequality holds since at least one of the two coefficients $2C$ and $S-2C$ is positive. This contradiction proves the lemma.
\end{proof}

\begin{theorem}
\label{height390:two-width-factor}
\label{height390:positive-first}\label{height390:zero-first}
Let $W$ be a nonempty closed height walk of span
at most two, starting at any height in its range. Then $++W$ has a perfect fractional
homomorphic tiling in every digraph $F$ with $\delta^0(F)>|F|/2$.
\end{theorem}
\begin{proof}
If $W$ has span one and starts at its lower height, then $W=(+-)^s$ for some $s\ge1$, and Theorem~\ref{imbalance402:isolated} applies. If it starts at its upper height, we reverse the
word and the edges of $F$. So we may assume that $W$ has span two. We write $E(j)=+(+-)^j-$, which is a closed walk from height zero to height zero visiting height two.

\emph{Starting at a boundary height.} Reversing if necessary, we may assume that $W$ starts at the bottom of its range, at height zero. Then $W=E(a_1)\cdots E(a_s)$ for some $a_i\ge0$, where $T=\sum_ia_i\ge1$, as $W$ has span two, and $|H|=2S$ with $S=s+T+1$.
If $|H|=6$, then $s=T=1$ and $++W=D_{4,2}$, which is covered by
Theorem~\ref{a:fractional}. Otherwise $|H|\ge8$.
The walk $W$ stays in $\{0,1,2\}$ and descends from two to zero, so the maximum drop of $H$ is two, and
\eqref{norm391:range} holds.

Delete the middle position of the initial pair $++$. The remaining path has height classes of multiplicities
$(s+1,s+T,T)$, from the lowest height upwards.
If $s\ge T-1$, the first two classes satisfy the conditions of Lemma~\ref{height390:heavy-pair}: their minimum
is at least $S/2$, and their sum minus $3S/2-1$ is
$(s-T+1)/2\ge0$.

Otherwise $s\le T-2$. Let $\ell\in\{0,\ldots,s-1\}$ be the number of
factors $E(0)=+-$ before the first factor $E(a_i)$ with $a_i\ge1$.
The first two positive edges of this factor form a new pair $++$, and we delete their middle position. In the height walk of $H$, this pair goes from height two to height four. Read the remaining path from height four, and then, after wrapping around, from the start of $H$ with heights translated by two.
The first part stays in $\{2,3,4\}$, and the second part consists of the old initial pair and $\ell$ factors $+-$, and stays in $\{2,\ldots,5\}$.
Counting positions gives class multiplicities
\[
 (s-\ell,\ T+s-\ell,\ T+\ell+1,\ \ell)
 \quad\text{at heights }2,3,4,5,
\]
where the last class is empty if $\ell=0$.
The middle two classes both have multiplicity at least $T+1\ge S/2$, and the sum of their multiplicities minus $3S/2-1$ is $(T-s+1)/2>0$.
Lemma~\ref{height390:heavy-pair} again applies.

\emph{Starting at the middle height.} Translate $W$ so that it starts at height one and stays in $\{0,1,2\}$,
and let $a,b\ge1$ be the numbers of steps from $0$ to $1$ and from $1$ to $2$, respectively, so that $S=a+b+1$. If $|H|=6$, then, depending on the order of the two excursions of $W$, the orientation $++W$ is, up to rotation,
$B$ or $D_{4,2}$, and we use
Theorem~\ref{base392:fractional}. So assume that $|H|\ge8$.
If the maximum drop of $H$ is one, then no two backward edges are consecutive, and
Theorem~\ref{imbalance402:isolated} applies. Thus we may assume that the maximum drop is two, so that \eqref{norm391:range} holds.

Deleting the middle position of the initial pair gives class multiplicities $(a,a+b+1,b)$.
If $a\ge b+1$, its first two classes satisfy the conditions of Lemma~\ref{height390:heavy-pair}. So suppose that $a\le b$. If $W$ starts with $+$,
we take the second initial positive edge and the first edge of $W$ as the new initial pair. The rest of the word is then a closed walk which starts at height two in $\{0,1,2\}$, and we are in the boundary case.
If $W$ ends with $+$, we take its last edge and the first initial positive edge; the rest is a closed walk starting at height zero, and again we are in the boundary case.

Otherwise, $W$ starts and ends with $-$, so it is a sequence of
lower excursions $-+$ and upper excursions $+-$.
Let $\ell\ge1$ be the number of lower excursions before the first upper one. If $\ell=a$,
then all lower excursions precede all upper ones, and then no two backward edges of $H$ are consecutive, contrary to the maximum drop being two. Thus $1\le\ell<a$.
Take as the new pair the two positive edges at the first transition from a lower excursion to an upper one, and delete their middle position. The remaining path has class multiplicities
\[
 (a-\ell,\ b+a-\ell,\ b+\ell+1,\ \ell)
 \quad\text{at heights }0,1,2,3,
\]
with both endpoints at height two. The middle two classes both have multiplicity
at least $b+1\ge S/2$, and the sum of their multiplicities minus $3S/2-1$ is
$(b-a+1)/2>0$. Lemma~\ref{height390:heavy-pair} applies once more.
\end{proof}
